\documentclass[10pt]{mypackage}

\usepackage[uprightgreeks,light]{kpfonts}
\renewcommand{\coloneq}{\coloneqq}
%\renewcommand{\mathbb}{\mathds}
\usepackage{homework}
%\usepackage{notes}

%\usepackage[ backend=bibtex, style = alphabetic, sorting=ynt ]{biblatex}
%\addbibresource{  }

\usepackage{parskip}

\fancyhf{}
\fancyhead[R]{Avinash Iyer}
\fancyhead[L]{Algebraic Topology: Homework 29}
\fancyfoot[C]{\thepage}

\setcounter{secnumdepth}{0}

\begin{document}
\RaggedRight
\begin{problem}[Problem 1]
  Prove that if $A,B\subseteq A\cup B$ are finite CW complexes, then
  \begin{align*}
    \chi\left( A\cup B \right) &= \chi(A) + \chi(B) - \chi\left( A\cap B \right).
  \end{align*}
\end{problem}
\begin{solution}
  By subdividing the cell complex where appropriate, we may assume that $A$, $B$, and $A\cap B$ are subcomplexes of $A\cup B$. By the principle of inclusion and exclusion, we have that
  \begin{align*}
    \operatorname{rk}\left( C_i(A\cup B) \right) &= \operatorname{rk}\left( C_i(A) \right) + \operatorname{rk}\left( C_i(B) \right) - \operatorname{rk}\left( C_i\left( A\cap B \right) \right),
  \end{align*}
  following from applying a counting argument to the $i$-cells of $A$, $B$, and $A\cap B$. Taking alternating sums at all levels gives
  \begin{align*}
    \chi\left( A\cup B \right) &= \sum_{i=0}^{n} \left( -1 \right)^{i} \operatorname{rk}\left( C_i\left( A\cup B \right) \right)\\
                               &= \sum_{i=0}^{n} \left( -1 \right)^{i} \left( \operatorname{rk}\left( A \right) + \operatorname{rk}\left( B \right) - \operatorname{rk}\left( A\cap B \right) \right)\\
                               &= \chi\left( A \right) + \chi\left( B \right) - \chi\left( A\cap B \right).
  \end{align*}
\end{solution}
\begin{problem}[Problem 2]
  Compute the homology groups of the space obtained from the torus $T^2 = S^1\times S^1$ by attaching a Möbius strip along its boundary to $S^1\times \set{x_0}\subseteq T^2$.
\end{problem}
\begin{solution}
  We may view the Möbius strip as a combination of one $0$-cell, one $1$-cell, and one $2$-cell, so that $\chi\left( B \right) = 1$. Similarly, we have $\chi\left( A \right) = 0$ since we know that the Euler characteristic of any orientable surface of genus $g$ is given by $2-2g$. Therefore, we have
  \begin{align*}
    \chi\left( A\cup B \right) &= \chi\left( A \right) + \chi\left( B \right) - \chi\left( A\cap B \right)\\
                               &= 1.
  \end{align*}

\end{solution}
\end{document}
